\documentclass[../../../include/open-logic-section]{subfiles}

\begin{document}

\olfileid{sth}{cardinals}{cp}
\olsection{د کانتور اصل}

خپل ذهن بېرته \olref[ordinals][vn]{sec} ته يوسئ۔ هلته مو د ښه
ترتيب شويو سټونو په اړه خبرې کولې او وړانديز مو وکړ چې ښه به وي
داسې څيزونه ولرو چې د ښه ترتيبونو استازيتوب وکړي۔ په همدې فکر مو
ترتيبي عددونه وپېژندل، او بيا مو په
\olref[ordinals][ordtype]{ordtypesworklikeyouwant} کښې وښودل چې دوی
هماغسې چلند کوي لکه چې غواړو، يعنې:
\[
	\ordtype{A, <} = \ordtype{B, \lessdot}
	\text{ هغه وخت او يوازې هغه وخت چې } \tuple{A, <} \isomorphic \tuple{B, \lessdot}.
\]
ذهن لا پخوا، \olref[sfr][siz][equ]{sec} ته يوسئ۔ هلته مو په ساده
بڼه د سټ د «اندازې» مفهوم وپېژانده۔ په ځانګړې توګه مو وويل چې دوه
سټونه، $\cardeq{A}{B}$، هغه وخت او يوازې هغه وخت هم‌شمېره دي چې
!!a{bijection} $f \colon A \to B$ موجود وي۔ دا مفهوم په خپل ذات
کښې د ښه ترتيب له مفهومه {ساده} دے: موږ يوازې له !!{bijection}s سره
کار لرو، او دا نه څېړو چې (لکه د ترتيب په هم‌شکلتيا کښې) دغه
!!{bijection}s «کوم جوړښت وساتي» که نه۔

له دې ټولو يو څرګند فکر راولاړېږي۔ لکه څنګه مو چې د ښه ترتيبونو د
سنجولو دپاره ځانګړي څيزونه، \emph{ترتيبي عددونه}، وپېژندل، همداسې
د اندازې د سنجولو دپاره هم ځانګړي څيزونه، \emph{اصلي عددونه}،
پېژندلے شو۔ د دې څپرکي موخه همدا ده۔

مخکښې له دې چې ووايو اصلي عددونه به څه وي، بايد هغه اصل وټاکو
چې دوی يې پوره کړي۔ که $\card{X}$ د سټ $X$ د اصلي شمېر دپاره
وليکو، غواړو چې لاندې شرط پوره کړي:
\[
	\card{A} = \card{B} \text{ هغه وخت او يوازې هغه وخت چې } \cardeq{A}{B}.
\]
دا به د \emph{کانتور} اصل وبولو، ځکه چې ښايي کانتور لومړنے کس و
چې دا يې په څرګنده توګه په ذهن کښې درلود۔ (له \emph{هيوم} اصل
سره د دې د اړيکې په اړه به په \olref[hp]{sec} کښې نور ووايو۔) نو
موخه مو دا ده چې $\card{X}$ د هر $X$ دپاره داسې تعريف کړو چې د
کانتور اصل ترې ترلاسه شي۔

\end{document}
